% Part: normal-modal-logic
% Chapter: completeness
% Section: lindenbaums-lemma

\documentclass[../../../include/open-logic-section]{subfiles}

\begin{document}

\olfileid{nml}{com}{lin}

\olsection{لم لیندنباوم}

لم لیندنباوم ثابت می‌کند هر مجموعهٔ~$\Sigma$-سازگار از !!{formula}ها
در دست‌کم یک مجموعهٔ~$\Sigma$-سازگار \emph{کامل} جای دارد. ساخت مدل
کانونی ما نشان خواهد داد که برای هر مجموعهٔ~$\Sigma$-سازگار کامل
$\Delta$، جهانی در مدل کانونی وجود دارد که همه و تنها !!{formula}های
$\Delta$ در آن صادق‌اند. پس لم لیندنباوم تضمین می‌کند هر مجموعهٔ
$\Sigma$-سازگار در جهانی از مدل کانونی صادق است.

\begin{thm}[لم لیندنباوم]\ollabel{thm:lindenbaum}
  اگر~$\Gamma$، $\Sigma$-سازگار باشد، مجموعهٔ~$\Sigma$-سازگار کامل
  $\Delta$ای وجود دارد که~$\Gamma$ را گسترش می‌دهد.
\end{thm}

\begin{proof}
بگذارید~$!A_0$، $!A_1$، \dots{} فهرستی فراگیر از همهٔ فرمول‌های زبان
باشد (تکرار مجاز است). برای نمونه، با فهرست‌کردن~$\Obj p_0$ آغاز کنید
و در هر مرحله~$n \ge 1$، فرمول‌های متناهی‌شمارِ با طول حداکثر~$n$ را
که تنها از متغیرهای~$\Obj p_0$، \dots،~$\Obj p_n$ استفاده می‌کنند
فهرست کنید. مجموعه‌های !!{formula}~$\Delta_n$ را با استقرا بر~$n$
تعریف می‌کنیم و سپس~$\Delta = \bigcup_n \Delta_n$ قرار می‌دهیم. نخست
$\Delta_0 = \Gamma$ قرار می‌دهیم. با فرض تعریف‌شدن~$\Delta_n$،
$\Delta_{n+1}$ را چنین تعریف می‌کنیم:
\[
\Delta_{n+1} =
\begin{cases}
  \Delta_n \cup \{!A_n\}, & \text{اگر~$\Delta_n \cup \{ !A_n\}$
    نسبت به~$\Sigma$ سازگار باشد؛} \\
  \Delta_n \cup \{ \lnot !A_n\}, & \text{در غیر این صورت.}
\end{cases}
\]
اکنون بگذارید~$\Delta = \bigcup_{n=0}^\infty \Delta_n$.

باید نشان دهیم این تعریف واقعاً مجموعهٔ~$\Delta$ای با خواص لازم به دست
می‌دهد؛ یعنی~$\Gamma \subseteq \Delta$ و~$\Delta$ کاملِ
$\Sigma$-سازگار است.

آشکار است~$\Gamma \subseteq \Delta$، زیرا بنا بر ساخت
$\Delta_0 \subseteq \Delta$ و~$\Delta_0 = \Gamma$. در واقع، برای هر
$n$، $\Delta_n \subseteq \Delta$، زیرا~$\Delta$ اجتماع همهٔ
$\Delta_n$هاست. (چون در هر گام ساخت !!a{formula} به مجموعهٔ ازپیش‌ساخته
می‌افزاییم، $\Delta_n \subseteq \Delta_{n+1}$؛ پس از تعدیِ~$\subseteq$
نتیجه می‌شود هرگاه~$n \le m$، $\Delta_n \subseteq \Delta_m$.) در هر
مرحلهٔ ساخت یا~$!A_n$ یا~$\lnot !A_n$ را می‌افزاییم و هر !!{formula}
دست‌کم یک بار در فهرست همهٔ~$!A_n$ها ظاهر می‌شود. پس برای هر~$!A$، یا
$!A \in \Delta$ یا~$\lnot !A \in \Delta$؛ بنابراین~$\Delta$ بنا به
تعریف کامل است.

سرانجام باید نشان دهیم~$\Delta$ نسبت به~$\Sigma$ سازگار است. برای این
کار نشان می‌دهیم (الف) اگر~$\Delta$ نسبت به~$\Sigma$ ناسازگار می‌بود،
یکی از~$\Delta_n$ها نسبت به~$\Sigma$ ناسازگار می‌بود، و (ب) همهٔ
$\Delta_n$ها نسبت به~$\Sigma$ سازگارند.

پس فرض کنید~$\Delta$ نسبت به~$\Sigma$ ناسازگار باشد. آنگاه
$\Delta \Proves[\Sigma] \lfalse$؛ یعنی~$!A_1$، \dots،~$!A_k \in
\Delta$ وجود دارند چنان‌که
$\Sigma \Proves !A_1 \lif (!A_2 \lif \cdots (!A_k \lif
\lfalse)\dots)$. چون~$\Delta = \bigcup_{n=0}^\infty \Delta_n$، هر
$!A_i$ برای یک~$n_i$ عضو~$\Delta_{n_i}$ است. بگذارید~$n$ بزرگ‌ترینِ
آن‌ها باشد. چون~$n_i \le n$، $\Delta_{n_i} \subseteq \Delta_n$. پس
همهٔ~$!A_i$ها عضو یک~$\Delta_n$ هستند. این بدان معناست که
$\Delta_n \Proves[\Sigma] \lfalse$؛ یعنی~$\Delta_n$ نسبت به~$\Sigma$
ناسازگار است.

برای نشان‌دادن اینکه هر~$\Delta_n$ نسبت به~$\Sigma$ سازگار است، از یک
استقرای ساده بر~$n$ استفاده می‌کنیم. $\Delta_0 = \Gamma$ و فرض کردیم
$\Gamma$ نسبت به~$\Sigma$ سازگار است. پس ادعا برای~$n=0$ برقرار است.
اکنون فرض کنید برای~$n$ برقرار باشد؛ یعنی~$\Delta_n$ نسبت به~$\Sigma$
سازگار باشد. اگر~$\Delta_n \cup \{!A_n\}$ نسبت به~$\Sigma$ سازگار
باشد، $\Delta_{n+1}$ همان است؛ در غیر این صورت
$\Delta_n \cup \{\lnot!A_n\}$ است. در حالت نخست~$\Delta_{n+1}$ آشکارا
نسبت به~$\Sigma$ سازگار است. اما بنا بر
\olref[prf][con]{prop:consistencyfacts}\olref[prf][con]{prop:consistencyfacts-c}،
یا~$\Delta_n \cup \{!A_n\}$ یا~$\Delta_n \cup \{\lnot!A_n\}$ سازگار
است؛ پس در حالت دیگر نیز~$\Delta_{n+1}$ سازگار است.
\end{proof}

\begin{cor}\ollabel{cor:provability-characterization}
  $\Gamma \Proves[\Sigma] !A$ اگر و تنها اگر~$!A \in \Delta$ برای هر
  مجموعهٔ~$\Sigma$-سازگار کامل~$\Delta$ که~$\Gamma$ را گسترش دهد
  (از جمله وقتی~$\Gamma = \emptyset$، که در این صورت توصیفی دیگر از
  دستگاه موجهاتی~$\Sigma$ به دست می‌آوریم).
\end{cor}

\begin{proof}
  فرض کنید~$\Gamma \Proves[\Sigma] !A$ و~$\Delta$ هر مجموعهٔ
  $\Sigma$-سازگار کامل گسترش‌دهندهٔ~$\Gamma$ باشد. اگر
  $!A \notin \Delta$، آنگاه بنا بر بیشینگی~$\lnot!A \in \Delta$؛ پس
  $\Delta \Proves[\Sigma] !A$ (بنا بر یکنوایی) و
  $\Delta \Proves[\Sigma] \lnot!A$ (بنا بر بازتابی‌بودن)، و بنابراین
  $\Delta$ ناسازگار است. برعکس، اگر~$\Gamma \Proves/[\Sigma] !A$،
  آنگاه~$\Gamma \cup \{\lnot!A\}$ نسبت به~$\Sigma$ سازگار است؛ و بنا
  بر لم لیندنباوم، مجموعهٔ سازگار کامل~$\Delta$ای وجود دارد که
  $\Gamma \cup \{\lnot!A\}$ را گسترش می‌دهد. بنا بر سازگاری،
  $!A \notin \Delta$.
\end{proof}

\end{document}
